% Options for packages loaded elsewhere
\PassOptionsToPackage{unicode}{hyperref}
\PassOptionsToPackage{hyphens}{url}
\PassOptionsToPackage{dvipsnames,svgnames,x11names}{xcolor}
\documentclass[
]{article}
\usepackage{xcolor}
\usepackage[margin=1in]{geometry}
\usepackage{amsmath,amssymb}
\setcounter{secnumdepth}{-\maxdimen} % remove section numbering
\usepackage{iftex}
\ifPDFTeX
  \usepackage[T1]{fontenc}
  \usepackage[utf8]{inputenc}
  \usepackage{textcomp} % provide euro and other symbols
\else % if luatex or xetex
  \usepackage{unicode-math} % this also loads fontspec
  \defaultfontfeatures{Scale=MatchLowercase}
  \defaultfontfeatures[\rmfamily]{Ligatures=TeX,Scale=1}
\fi
\usepackage{lmodern}
\ifPDFTeX\else
  % xetex/luatex font selection
\fi
% Use upquote if available, for straight quotes in verbatim environments
\IfFileExists{upquote.sty}{\usepackage{upquote}}{}
\IfFileExists{microtype.sty}{% use microtype if available
  \usepackage[]{microtype}
  \UseMicrotypeSet[protrusion]{basicmath} % disable protrusion for tt fonts
}{}
\makeatletter
\@ifundefined{KOMAClassName}{% if non-KOMA class
  \IfFileExists{parskip.sty}{%
    \usepackage{parskip}
  }{% else
    \setlength{\parindent}{0pt}
    \setlength{\parskip}{6pt plus 2pt minus 1pt}}
}{% if KOMA class
  \KOMAoptions{parskip=half}}
\makeatother
\usepackage{longtable,booktabs,array}
\newcounter{none} % for unnumbered tables
\usepackage{calc} % for calculating minipage widths
% Correct order of tables after \paragraph or \subparagraph
\usepackage{etoolbox}
\makeatletter
\patchcmd\longtable{\par}{\if@noskipsec\mbox{}\fi\par}{}{}
\makeatother
% Allow footnotes in longtable head/foot
\IfFileExists{footnotehyper.sty}{\usepackage{footnotehyper}}{\usepackage{footnote}}
\makesavenoteenv{longtable}
\setlength{\emergencystretch}{3em} % prevent overfull lines
\providecommand{\tightlist}{%
  \setlength{\itemsep}{0pt}\setlength{\parskip}{0pt}}
% pdf_header.tex — Unicode-to-LaTeX mappings for pandoc/xelatex builds
% Used by build_pdfs.sh for ManyClocks.pdf and Equations.pdf.
% Maps prose-mode Unicode characters (Greek letters, math symbols, arrows,
% blackboard bold, etc.) to LaTeX commands so they render correctly in
% paragraphs, captions, and tables — outside of $...$ math mode where
% the math font would normally handle them.

\usepackage{newunicodechar}
\usepackage{amsmath,amssymb}
\newunicodechar{ℏ}{\ensuremath{\hbar}}
\newunicodechar{⟨}{\ensuremath{\langle}}
\newunicodechar{⟩}{\ensuremath{\rangle}}
\newunicodechar{→}{\ensuremath{\rightarrow}}
\newunicodechar{←}{\ensuremath{\leftarrow}}
\newunicodechar{↔}{\ensuremath{\leftrightarrow}}
\newunicodechar{⇒}{\ensuremath{\Rightarrow}}
\newunicodechar{⇔}{\ensuremath{\Leftrightarrow}}
\newunicodechar{≈}{\ensuremath{\approx}}
\newunicodechar{≠}{\ensuremath{\neq}}
\newunicodechar{≤}{\ensuremath{\leq}}
\newunicodechar{≥}{\ensuremath{\geq}}
\newunicodechar{≪}{\ensuremath{\ll}}
\newunicodechar{≫}{\ensuremath{\gg}}
\newunicodechar{≲}{\ensuremath{\lesssim}}
\newunicodechar{≳}{\ensuremath{\gtrsim}}
\newunicodechar{∝}{\ensuremath{\propto}}
\newunicodechar{∼}{\ensuremath{\sim}}
\newunicodechar{√}{\ensuremath{\sqrt{}}}
\newunicodechar{∞}{\ensuremath{\infty}}
\newunicodechar{∂}{\ensuremath{\partial}}
\newunicodechar{∇}{\ensuremath{\nabla}}
\newunicodechar{∫}{\ensuremath{\int}}
\newunicodechar{∑}{\ensuremath{\sum}}
\newunicodechar{∈}{\ensuremath{\in}}
\newunicodechar{∉}{\ensuremath{\notin}}
\newunicodechar{⊗}{\ensuremath{\otimes}}
\newunicodechar{⊕}{\ensuremath{\oplus}}
\newunicodechar{·}{\ensuremath{\cdot}}
\newunicodechar{×}{\ensuremath{\times}}
\newunicodechar{±}{\ensuremath{\pm}}
\newunicodechar{✓}{\ensuremath{\checkmark}}
\newunicodechar{ψ}{\ensuremath{\psi}}
\newunicodechar{Ψ}{\ensuremath{\Psi}}
\newunicodechar{φ}{\ensuremath{\varphi}}
\newunicodechar{Φ}{\ensuremath{\Phi}}
\newunicodechar{χ}{\ensuremath{\chi}}
\newunicodechar{σ}{\ensuremath{\sigma}}
\newunicodechar{Σ}{\ensuremath{\Sigma}}
\newunicodechar{α}{\ensuremath{\alpha}}
\newunicodechar{β}{\ensuremath{\beta}}
\newunicodechar{γ}{\ensuremath{\gamma}}
\newunicodechar{δ}{\ensuremath{\delta}}
\newunicodechar{Δ}{\ensuremath{\Delta}}
\newunicodechar{ε}{\ensuremath{\varepsilon}}
\newunicodechar{ω}{\ensuremath{\omega}}
\newunicodechar{Ω}{\ensuremath{\Omega}}
\newunicodechar{η}{\ensuremath{\eta}}
\newunicodechar{θ}{\ensuremath{\theta}}
\newunicodechar{Θ}{\ensuremath{\Theta}}
\newunicodechar{ν}{\ensuremath{\nu}}
\newunicodechar{ρ}{\ensuremath{\rho}}
\newunicodechar{τ}{\ensuremath{\tau}}
\newunicodechar{λ}{\ensuremath{\lambda}}
\newunicodechar{Λ}{\ensuremath{\Lambda}}
\newunicodechar{μ}{\ensuremath{\mu}}
\newunicodechar{π}{\ensuremath{\pi}}
\newunicodechar{Π}{\ensuremath{\Pi}}
\newunicodechar{ξ}{\ensuremath{\xi}}
\newunicodechar{ζ}{\ensuremath{\zeta}}
\newunicodechar{κ}{\ensuremath{\kappa}}
\newunicodechar{ℓ}{\ensuremath{\ell}}
\newunicodechar{ℝ}{\ensuremath{\mathbb{R}}}
\newunicodechar{ℂ}{\ensuremath{\mathbb{C}}}
\newunicodechar{ℕ}{\ensuremath{\mathbb{N}}}
\newunicodechar{ℤ}{\ensuremath{\mathbb{Z}}}
\newunicodechar{²}{\textsuperscript{2}}
\newunicodechar{³}{\textsuperscript{3}}
\newunicodechar{¹}{\textsuperscript{1}}
\newunicodechar{₀}{\textsubscript{0}}
\newunicodechar{₁}{\textsubscript{1}}
\newunicodechar{₂}{\textsubscript{2}}
\newunicodechar{⁻}{\textsuperscript{-}}
\newunicodechar{⁰}{\textsuperscript{0}}
\newunicodechar{⁴}{\textsuperscript{4}}
\newunicodechar{⁵}{\textsuperscript{5}}
\newunicodechar{⁶}{\textsuperscript{6}}
\newunicodechar{⁷}{\textsuperscript{7}}
\newunicodechar{⁸}{\textsuperscript{8}}
\newunicodechar{⁹}{\textsuperscript{9}}
\newunicodechar{ᵃ}{\textsuperscript{a}}
\newunicodechar{ᵇ}{\textsuperscript{b}}
\newunicodechar{ᵀ}{\textsuperscript{T}}
\newunicodechar{₇}{\textsubscript{7}}
\newunicodechar{⟹}{\ensuremath{\Longrightarrow}}
\newunicodechar{−}{$-$}
\newunicodechar{—}{---}
\newunicodechar{–}{--}
\newunicodechar{…}{\ldots}
\newunicodechar{′}{\ensuremath{{}^{\prime}}}
\usepackage{bookmark}
\IfFileExists{xurl.sty}{\usepackage{xurl}}{} % add URL line breaks if available
\urlstyle{same}
\hypersetup{
  pdftitle={Cosmic Expansion as Surface-Area Particle Creation in a Dirac--Kuramoto Substrate},
  pdfauthor={John Bramble, MD},
  colorlinks=true,
  linkcolor={Maroon},
  filecolor={Maroon},
  citecolor={Blue},
  urlcolor={Blue},
  pdfcreator={LaTeX via pandoc}}

\title{Cosmic Expansion as Surface-Area Particle Creation in a
Dirac--Kuramoto Substrate}
\usepackage{etoolbox}
\makeatletter
\providecommand{\subtitle}[1]{% add subtitle to \maketitle
  \apptocmd{\@title}{\par {\large #1 \par}}{}{}
}
\makeatother
\subtitle{A Coherence-Based Reframing of the FRW Universe and the
Cosmological Constant}
\author{John Bramble, MD}
\date{May 2026}

\begin{document}
\maketitle

{
\setcounter{tocdepth}{2}
\tableofcontents
}
\textbf{John Bramble, MD¹}

¹ Independent Research

\emph{Correspondence: rayolddog (GitHub)}

\emph{Preprint --- not yet peer reviewed}

\emph{AI Disclosure: This work was developed in collaboration with
Claude Opus 4.7 (Anthropic), as described in the Author Contributions
section. Per current journal guidelines, LLMs do not satisfy authorship
criteria; the human author bears full responsibility for all content.}

\emph{Companion paper to: Bramble, J. (2026), ``Many Clocks
Interpretation of Quantum Mechanics: Mass as Chiral Coupling,
Re-synchronization in the Bulk as Measurement'' {[}PAPER\_UNIFIED, under
review{]}.}

\begin{center}\rule{0.5\linewidth}{0.5pt}\end{center}

\subsection{Abstract}\label{abstract}

The standard Friedmann--Robertson--Walker (FRW) cosmology treats
expansion as a geometric property of spacetime sourced by a cosmological
constant Λ inserted by hand into Einstein's field equations. Λ
originated as Einstein's 1917 patch to recover a \emph{static} universe
--- a Newtonian-era prejudice --- and was later resurrected as ``dark
energy'' to fit accelerating expansion. In both incarnations, Λ enters
the framework axiomatically rather than emergently.

This paper proposes a reframing in which cosmic expansion is \emph{not}
a property of an underlying expanding spacetime, but instead the
dynamical growth of a finite \emph{populated region} embedded in an
infinite, static quantum-field substrate. Expansion proceeds via
real-particle creation from vacuum fluctuations in regions of low
temperature and low gravitational density --- the regime in which the
Dirac--Kuramoto (DK) decoherence rates Γ\_env(T) and Γ\_grav(M) both
vanish. Because creation occurs preferentially at the boundary of the
populated region, the rate of expansion is proportional to the
\textbf{surface area} of the universe rather than its volume, naturally
producing a Hubble rate that scales as H ∝ 1/R and an effective
cosmological term Λ\_eff ∝ 1/R² --- matching the observed Λ
\textasciitilde{} 1/L\_H² without postulation.

The framework predicts three observational signatures distinct from
ΛCDM: (1) the local Hubble rate should anti-correlate with local matter
density, yielding the observed H₀ tension between CMB and
local-distance-ladder measurements; (2) cosmic voids should expand
faster than virialized structures by a measurable margin; (3) the
surface-area scaling implies an ``emergent FRW metric'' inside the
populated region that preserves Tolman (1+z)⁴ dimming and CMB blackbody
structure, while the underlying substrate remains Minkowski. The
proposal is presented as a working-paper sketch, not a fully derived
theory; the load-bearing assumption --- that excitations of the
bulk-coherence order parameter Φ\_bulk inherit equation of state w = −1
through their Goldstone-mode character --- is identified as the next
required derivation.

\begin{center}\rule{0.5\linewidth}{0.5pt}\end{center}

\subsection{1. Introduction}\label{introduction}

\subsubsection{1.1 The framework problem with
Λ}\label{the-framework-problem-with-ux3bb}

Modern cosmology rests on the Friedmann--Robertson--Walker (FRW) metric,
an exact solution of Einstein's field equations under the
\emph{Cosmological Principle} --- the assumption that the universe is
homogeneous and isotropic at large scales. The Friedmann equations
governing the time evolution of the FRW scale factor a(t) take the form:

\[H^2 \equiv \left(\frac{\dot a}{a}\right)^2 = \frac{8 \pi G}{3} \rho - \frac{k c^2}{a^2} + \frac{\Lambda c^2}{3} \qquad (\text{Friedmann I})\]

\[\frac{\ddot a}{a} = -\frac{4 \pi G}{3}\left(\rho + \frac{3 p}{c^2}\right) + \frac{\Lambda c^2}{3} \qquad (\text{Friedmann II})\]

where ρ is the total energy density, p its pressure, k ∈ \{−1, 0, +1\}
the spatial-curvature index, and Λ the cosmological constant. The
framework is extraordinarily successful: it reproduces Hubble's law
(recession velocity proportional to distance), predicts the cosmic
microwave background (CMB) and its near-perfect blackbody spectrum,
accounts for primordial nucleosynthesis, and --- with the addition of a
positive Λ --- fits the accelerating expansion measured by Type Ia
supernovae from 1998 onward.

The framework's success does not, however, dissolve the conceptual
problem with Λ itself. Λ enters Einstein's equations as a free parameter
chosen to match observation. Its origin is unspecified by the theory.
Three distinct unsolved problems attach to it:

\begin{enumerate}
\def\labelenumi{\arabic{enumi}.}
\item
  \textbf{The cosmological-constant problem proper.} Quantum field
  theory predicts a vacuum energy density of order
  \(\rho_{\text{vac}} \sim M_{\text{Pl}}^4\), roughly 120 orders of
  magnitude larger than the observed
  \(\rho_\Lambda \sim 10^{-47}\ \text{GeV}^4\). Why these two quantities
  differ so radically is the largest known discrepancy between theory
  and observation in all of physics.
\item
  \textbf{The coincidence problem.} Why is \(\rho_\Lambda\) comparable
  in magnitude to \(\rho_{\text{matter}}\) in the current cosmic epoch,
  given that the two have entirely different scaling with \(a(t)\)?
\item
  \textbf{The H₀ tension.} Independent measurements of the current
  expansion rate disagree at \textasciitilde5σ. CMB-anchored
  measurements (Planck) give \(H_0 \approx 67.4\ \text{km/s/Mpc}\);
  local distance-ladder measurements (SH0ES) give
  \(H_0 \approx 73.0\ \text{km/s/Mpc}\). ΛCDM with a single, uniform Λ
  cannot accommodate both.
\end{enumerate}

Each of these problems suggests that Λ is \emph{not} the right primitive
--- that expansion ought to emerge from some deeper dynamics, not be
postulated.

\subsubsection{1.2 The Dirac--Kuramoto framework: what is already
available}\label{the-dirackuramoto-framework-what-is-already-available}

The DK framework {[}companion paper, PAPER\_UNIFIED{]} proposes that
quantum measurement consists of two stages: (Stage 1)
pair-synchronization of a perturbed partner with a macroscopic
\emph{bulk} characterized by an order-parameter phase
\(\Phi_{\text{bulk}}\), followed by (Stage 2) irreversible
re-equilibration of the partner to the bulk via two parallel rate
channels,

\[\Gamma_{\text{Stage 2}}^{\text{total}} = \Gamma_{\text{env}}(T,\, J(\omega)) + \Gamma_{\text{grav}}(M,\, \Delta z) \qquad (3.5')\]

where \(\Gamma_{\text{env}}\) is the environmental (thermal +
bath-mediated) relaxation rate and \(\Gamma_{\text{grav}}\) is the
gravitational self-relocking rate of the bulk. Stage-2 events are the
rare bath-coupling moments that produce classical, irreversible
measurement records.

Both rates vanish in the limit of low temperature and low gravitational
density: \(\Gamma_{\text{env}} \to J(\omega)\) as \(T \to 0\) (the
spontaneous-emission floor), and \(\Gamma_{\text{grav}} \to 0\) as the
local matter density goes to zero. The DK framework therefore already
specifies the conditions under which bulk-coherence-mediated decoherence
ceases to act.

The proposal of this paper is that \emph{exactly those conditions} ---
low T and low gravity --- define the regime in which the cosmological
substrate converts vacuum fluctuations into real particles, and that
this conversion is the physical source of cosmic expansion.

\begin{center}\rule{0.5\linewidth}{0.5pt}\end{center}

\subsection{2. A Brief History of FRW
Cosmology}\label{a-brief-history-of-frw-cosmology}

\subsubsection{2.1 Einstein's static universe
(1917)}\label{einsteins-static-universe-1917}

Einstein applied his newly completed field equations of general
relativity to the universe as a whole and found that no static solution
existed for a homogeneous matter distribution under attractive gravity
alone --- the universe would either expand or contract. Believing on
philosophical grounds that the universe must be eternal and static,
Einstein modified his equations by adding a term \(\Lambda g_{\mu\nu}\),

\[G_{\mu\nu} + \Lambda g_{\mu\nu} = \frac{8 \pi G}{c^4} T_{\mu\nu},\]

where \(\Lambda\) was tuned to exactly balance gravitational attraction
at cosmic scales. The resulting ``Einstein static universe'' was a
closed 3-sphere of fixed radius.

\subsubsection{2.2 Friedmann and Lemaître (1922,
1927)}\label{friedmann-and-lemauxeetre-1922-1927}

The Russian mathematician Alexander Friedmann (1922), and independently
the Belgian priest-physicist Georges Lemaître (1927), solved Einstein's
equations under the assumptions of homogeneity and isotropy
\emph{without} demanding stasis. They found that the universe must in
general expand or contract, with the dynamics governed by Friedmann I
(above). Lemaître went further, connecting the expanding solutions to
observational redshift data and proposing what would later be called the
Big Bang. Einstein famously rejected Lemaître's physics --- ``your
calculations are correct, but your physics is abominable'' --- before
reversing himself after Hubble's 1929 observations.

\subsubsection{2.3 Hubble's law (1929)}\label{hubbles-law-1929}

Edwin Hubble's observation that galactic recession velocity is
proportional to distance,

\[v = H_0 \, d \qquad (\text{Hubble's law})\]

established expansion empirically. Einstein subsequently called his
introduction of \(\Lambda\) his ``greatest blunder.''

\subsubsection{2.4 Robertson and Walker (1935--36); FRW
metric}\label{robertson-and-walker-193536-frw-metric}

Howard Robertson and Arthur Walker independently proved that the most
general metric consistent with homogeneity and isotropy takes the form

\[ds^2 = -c^2 dt^2 + a(t)^2 \left[ \frac{dr^2}{1 - k r^2} + r^2 \left(d\theta^2 + \sin^2\theta\, d\phi^2 \right) \right] \qquad (\text{FRW metric})\]

where \(a(t)\) is the scale factor encoding all time evolution and
\(k \in \{-1, 0, +1\}\) labels open, flat, and closed spatial geometries
respectively. Observations of CMB power-spectrum peaks indicate
\(k = 0\) to within a fraction of a percent.

\subsubsection{2.5 The accelerating universe (1998) and Λ's
return}\label{the-accelerating-universe-1998-and-ux3bbs-return}

Type Ia supernova surveys (Riess et al.~1998; Perlmutter et al.~1999)
showed that the universe's expansion is \emph{accelerating}. Within FRW,
only a component with sufficiently negative pressure --- equation of
state \(w < -1/3\) --- can drive acceleration. The simplest such
component is a positive cosmological constant (\(w = -1\)), and ``ΛCDM''
became the standard cosmological model. The dark-energy term thereby
resurrected was mathematically Einstein's original \(\Lambda\) but with
no static-universe motivation --- purely an empirical fit.

\subsubsection{2.6 The standing question}\label{the-standing-question}

In its modern form, Λ is a single number tuned to one observation
(acceleration) that fits a wide range of others (CMB, large-scale
structure, BAO) --- but with the three open problems noted in §1.1. The
framework proposed below leaves the FRW machinery intact as an
\emph{effective} description of expansion inside the populated region,
while replacing the fundamental ontology with one in which Λ is
emergent.

\begin{center}\rule{0.5\linewidth}{0.5pt}\end{center}

\subsection{3. The Substrate--Population
Distinction}\label{the-substratepopulation-distinction}

\subsubsection{3.1 The standard fusion}\label{the-standard-fusion}

In standard FRW cosmology, two distinct ontological commitments are
fused:

\begin{itemize}
\tightlist
\item
  \textbf{Spacetime is the universe.} The FRW metric describes the
  entirety of physical space; there is no ``outside'' of the universe in
  the framework.
\item
  \textbf{The metric itself expands.} The scale factor \(a(t)\)
  describes expansion of space, not just motion of matter through space.
\end{itemize}

These commitments are taken together as a single package, and the
cosmological-constant problem inherits both: \(\Lambda\) is both the
energy density of spacetime itself and the agent of its expansion.

\subsubsection{3.2 The DK ontology}\label{the-dk-ontology}

The DK framework suggests a separation:

\begin{itemize}
\item
  \textbf{The substrate}: an infinite, eternal quantum-field structure
  (the vacuum and its fluctuations). The substrate has no privileged
  scale, no boundary, and is not itself expanding. Its metric, in the
  absence of coherent excitation, is Minkowski.
\item
  \textbf{The populated region}: the finite, growing subset of the
  substrate in which real-particle excitations exist and are
  bulk-coherent in the DK sense. This is what we call ``the universe.''
  It has a boundary, a scale, and a dynamical evolution.
\end{itemize}

Cosmic expansion, on this view, is \textbf{not} the metric of the
substrate stretching. It is the \emph{boundary of the populated region
advancing into the substrate} as vacuum fluctuations cross a coherence
threshold and become real excitations.

\subsubsection{3.3 Why the substrate metric remains
Minkowski}\label{why-the-substrate-metric-remains-minkowski}

The substrate is the unexcited vacuum of quantum field theory. QFT
vacuum states are Lorentz invariant --- they look the same in every
inertial frame. The unique Lorentz-invariant metric is Minkowski. So the
substrate, by its Lorentz symmetry alone, has no expansion to speak of:
there is no scale factor in Minkowski space because there is no scale at
all.

What we observe as ``expansion'' must therefore be a property of the
populated region's internal structure, not of the substrate.

\subsubsection{3.4 The effective FRW metric inside the populated
region}\label{the-effective-frw-metric-inside-the-populated-region}

Within the populated region, \(\Phi_{\text{bulk}}\) is established and
matter-energy is present at non-zero density. The DK framework already
predicts that bulk coherence determines an effective metric on the
coherent region: the order-parameter phase \(\Phi_{\text{bulk}}\)
couples to local proper-time evolution and to gravitational
pair-locking. Internally, observers measure an effective metric that ---
by symmetry, since the populated region is homogeneous and isotropic on
average --- takes FRW form. The scale factor \(a(t)\) that appears in
this effective metric is determined by the dynamics of the populated
region itself, not by any property of the substrate.

This is structurally analogous to \emph{analogue gravity} in
condensed-matter systems: phonons propagating in a Bose--Einstein
condensate experience an effective Lorentzian metric set by the
condensate density and flow {[}Unruh 1981, Barceló--Liberati--Visser
2005{]}. The analogue metric is real for its internal excitations even
though the underlying atoms live in flat lab space. Likewise, the
effective FRW metric is real for matter inside the cosmological
populated region, even though the underlying QFT substrate is Minkowski.

\begin{center}\rule{0.5\linewidth}{0.5pt}\end{center}

\subsection{4. The Creation Mechanism: Low Γ as the
Threshold}\label{the-creation-mechanism-low-ux3b3-as-the-threshold}

\subsubsection{4.1 Vacuum fluctuations and the virtual--real
distinction}\label{vacuum-fluctuations-and-the-virtualreal-distinction}

Quantum field theory predicts continuous vacuum fluctuations: pairs of
virtual particles appearing and annihilating on timescales bounded by
the energy--time uncertainty relation
\(\Delta E \cdot \Delta t \sim \hbar\). Under ordinary conditions,
virtual pairs annihilate before they can be detected or do measurable
work.

Several known mechanisms convert virtual pairs into \emph{real}
particles by preventing the annihilation:

{\def\LTcaptype{none} % do not increment counter
\begin{longtable}[]{@{}lll@{}}
\toprule\noalign{}
Mechanism & Trigger & Conversion source \\
\midrule\noalign{}
\endhead
\bottomrule\noalign{}
\endlastfoot
Schwinger pair production & Strong electric field & Field energy \\
Hawking radiation & Event horizon & Horizon geometry \\
Unruh effect & Acceleration & Kinetic energy of accelerated observer \\
Parker cosmological creation & Expanding metric & Metric work \\
Dynamical Casimir effect & Moving boundary & Boundary kinetic energy \\
\end{longtable}
}

Each of these mechanisms identifies a \emph{separation} of the virtual
pair before they recombine. The energy that becomes the real particles'
rest-mass-equivalent is supplied by the external agent (field, horizon,
acceleration, expansion, boundary motion).

\subsubsection{4.2 The DK threshold: low Γ as a fifth conversion
mechanism}\label{the-dk-threshold-low-ux3b3-as-a-fifth-conversion-mechanism}

The DK framework specifies a different conversion mechanism. In ordinary
matter, virtual pairs are re-absorbed because the bulk's coherence
dynamics --- specifically \(\Gamma_{\text{env}}(T)\) and
\(\Gamma_{\text{grav}}(M, \Delta z)\) --- quickly damp any fluctuation
back into the bulk-coherent state before it can become a
measurement-relevant real excitation. Bulk coherence acts as an active
\emph{suppressor} of vacuum fluctuations.

In regions of low temperature and low gravitational density, both rates
vanish:

\[\Gamma_{\text{env}}(T) \to J(\omega) \cdot [1 + 2 n_{\text{th}}(\omega, T)] \xrightarrow{T \to 0} J(\omega)\]

\[\Gamma_{\text{grav}}(M, \Delta z) \to 0 \quad \text{as} \quad M_{\text{local}} \to 0\]

When both rates fall below the inverse lifetime of a virtual
fluctuation, the fluctuation is no longer re-absorbed in time. It
survives long enough to become a real excitation of the field --- a real
particle. The suppression mechanism has failed.

\textbf{The proposal:} \emph{cosmic expansion is sourced by the failure
of bulk- coherence suppression in low-Γ regions of the substrate.} The
same vacuum fluctuations that occur ubiquitously become real only where
Γ\_env(T) and Γ\_grav(M) cannot damp them in time. These regions are:
(i) cosmic voids, (ii) the cold dark sky between galaxies and clusters,
and (iii) most critically, the \emph{boundary} of the populated region
itself, where the density of pre-existing matter is by definition
lowest.

\subsubsection{4.3 Where the energy comes
from}\label{where-the-energy-comes-from}

The energy budget for new particles must come from somewhere. The DK
framework points to a natural source: the \textbf{incompleteness of bulk
coherence at horizon scales}. Φ\_bulk cannot complete its sync beyond
the causal horizon; the residual coherence deficit is an energy
reservoir analogous to the spontaneous-emission floor \(J(\omega)\) in
the Stage-2 environmental term. When a virtual pair near the horizon
fails to be reabsorbed, it draws on this reservoir to upgrade to real-
particle status. The expansion of the populated region thereby converts
horizon-incomplete-coherence energy into new real-particle excitations.

\subsubsection{4.4 The equation-of-state
problem}\label{the-equation-of-state-problem}

A standing technical problem remains: \emph{real particles created from
the vacuum carry positive energy density, which gravitates, producing
contraction rather than expansion}. For the mechanism to drive
expansion, the created excitations must inherit an equation of state
\(w \approx -1\) (the cosmological-constant equation of state) rather
than \(w = 0\) (cold matter) or \(w = 1/3\) (radiation).

The DK framework provides a natural candidate: the created excitations
are not localized matter quanta but \textbf{collective phase modes
(Goldstone modes) of the broken-U(1) order parameter Φ\_bulk}. Goldstone
modes of a condensate are massless by Goldstone's theorem, propagate
relativistically in the absence of a preferred frame, and --- in
Lorentz-covariant formulations --- carry vacuum-like equation of state
\(w = -1\). The created excitations are thereby identified with
\emph{expansion-driving} quanta, not gravitating matter.

This identification is the load-bearing assumption of the proposal and
the next step in its derivation. A complete treatment must show, from
the broken-U(1) structure of cosmological Φ\_bulk, that its Goldstone
modes carry \(w = -1\) to the required precision and that their creation
rate at low Γ reproduces the observed expansion history.

\begin{center}\rule{0.5\linewidth}{0.5pt}\end{center}

\subsection{5. The Surface-Area Scaling}\label{the-surface-area-scaling}

\subsubsection{5.1 The geometric
observation}\label{the-geometric-observation}

Real-particle creation occurs preferentially where Γ\_env and Γ\_grav
are smallest. Within the populated region, both rates are non-zero
(matter is present, and the CMB sets a non-zero temperature floor). At
the \emph{boundary} of the populated region, however, both rates vanish
--- by definition, this is where the matter density of the populated
region goes to zero.

The boundary is approximately a 2-sphere of radius \(R(t)\) (set by the
extent of the populated region in cosmic time), with surface area

\[A(t) = 4 \pi R(t)^2.\]

The rate of real-particle creation at the boundary is proportional to
the area of the boundary times a per-area creation rate
\(\sigma_{\text{c}}\) set by the strength of vacuum fluctuations and the
relevant coupling:

\[\frac{dN}{dt} = \sigma_{\text{c}} \cdot A(t) = 4 \pi \sigma_{\text{c}} R(t)^2 \qquad (5.1)\]

\subsubsection{5.2 The Hubble rate as a function of
R}\label{the-hubble-rate-as-a-function-of-r}

If the rate of \emph{volume} growth of the populated region is
proportional to its surface area times a characteristic ``creation-front
speed'' \(v_{\text{c}}\),

\[\frac{dV}{dt} = v_{\text{c}} \cdot A(t),\]

then for a spherical region \(V = \frac{4}{3} \pi R^3\), we have
\(\frac{dV}{dt} = 4 \pi R^2 \dot R\), giving

\[\dot R = v_{\text{c}} \qquad (5.2)\]

--- the \emph{boundary advances at constant speed} \(v_{\text{c}}\).
Identifying the scale factor \(a(t) \propto R(t)\), the effective Hubble
rate is

\[H = \frac{\dot a}{a} = \frac{\dot R}{R} = \frac{v_{\text{c}}}{R(t)} \qquad (5.3)\]

The Hubble rate naturally \textbf{decreases as \(1/R\)} as the universe
grows. At fixed \(v_{\text{c}}\), an older, larger universe expands more
slowly in fractional terms even as its boundary advances at constant
absolute speed.

\subsubsection{5.3 The effective cosmological
constant}\label{the-effective-cosmological-constant}

Comparing (5.3) with the Λ-dominated form of Friedmann I,

\[H^2 = \frac{\Lambda c^2}{3} \implies \Lambda_{\text{eff}} = \frac{3 H^2}{c^2} = \frac{3 v_{\text{c}}^2}{c^2 R^2} \qquad (5.4)\]

so the effective cosmological constant scales as \textbf{1/R²}, matching
the observed scaling \(\Lambda \sim 1/L_H^2\) where
\(L_H \approx c/H_0\) is the Hubble length.

This is, at the order-of-magnitude level, the correct value for the
observed Λ --- without postulating Λ at all. The framework produces a
dynamically generated cosmological term whose magnitude is set by the
size of the populated region.

\subsubsection{5.4 Connection to the holographic
principle}\label{connection-to-the-holographic-principle}

The surface-area scaling has a deep precedent in physics: the
\textbf{holographic principle} (Bekenstein, 't Hooft, Susskind,
Verlinde), which holds that the maximum entropy enclosable in a region
is bounded not by its volume but by its surface area in Planck units,

\[S \leq \frac{A}{4 \ell_P^2}.\]

Verlinde's \emph{entropic gravity} and related proposals extend the
holographic principle to suggest that gravity itself, and possibly the
cosmological constant, arise from entropy gradients at boundaries.

The DK proposal is in the same family: expansion is sourced at the
horizon, with rate proportional to horizon area, and the effective Λ is
set by the inverse-area scale of the horizon. The DK framework differs
from entropic gravity in providing a \emph{physical mechanism} (low-Γ
vacuum-fluctuation conversion) rather than a purely thermodynamic
argument --- but the geometric scaling is the same.

\begin{center}\rule{0.5\linewidth}{0.5pt}\end{center}

\subsection{6. Testable Predictions}\label{testable-predictions}

\subsubsection{6.1 The H₀ tension}\label{the-hux2080-tension}

In standard ΛCDM, the present-day Hubble rate \(H_0\) is a single number
characterizing a uniform expansion. The measured tension between
CMB-derived (Planck: \(H_0 \approx 67.4\) km/s/Mpc) and
local-distance-ladder (SH0ES: \(H_0 \approx 73.0\) km/s/Mpc) values
represents a \textasciitilde5σ disagreement that ΛCDM cannot accommodate
without modification.

The DK framework predicts that \(H_{\text{local}}\) should depend on the
local matter density along the line of sight. The local distance ladder
samples regions inside our cosmic neighborhood, including the local void
(``KBC void'') in which the Milky Way may sit. CMB measurements average
over the entire observable horizon, sampling cosmic-mean density.

If expansion is sourced preferentially in low-density regions, then a
local probe biased toward void regions will measure a \emph{higher}
effective \(H_0\) than a global probe --- qualitatively matching the
observed tension.

\textbf{Quantitative prediction}: \(H_{\text{local}}/H_{\text{global}}\)
should correlate with the ratio of local-mean-density to
cosmic-mean-density, with the correlation function determined by the
relative weights of \(\Gamma_{\text{env}}\) and \(\Gamma_{\text{grav}}\)
suppression in the two density regimes. A first-order estimate predicts
a tension of order
\(\Delta H_0 / H_0 \sim (\rho_{\text{cosmic}} - \rho_{\text{local}}) / \rho_{\text{cosmic}}\);
the observed tension of \(\Delta H_0 / H_0 \approx 0.08\) is consistent
with the void underdensities of order \textasciitilde10\% inferred from
galaxy surveys.

\subsubsection{6.2 Void expansion rates}\label{void-expansion-rates}

ΛCDM predicts that cosmic voids expand at the same Hubble rate as the
surrounding universe, modified only by the local gravitational deficit
(which slightly \emph{enhances} recession). The DK framework predicts an
\emph{additional} contribution: voids should expand faster because real-
particle creation is actively occurring within them, sourcing additional
expansion not present in dense regions.

\textbf{Quantitative prediction}: the void expansion velocity excess
over ΛCDM should scale with the void's emptiness --- deeper voids expand
faster, with the excess vanishing as the void's mean density approaches
the cosmic mean.

This is in principle measurable through redshift-space distortions in
galaxy surveys around well-characterized voids (BOSS, eBOSS, DESI). A
positive detection would constitute direct evidence for the proposed
mechanism; a null result at the predicted level would falsify it.

\subsubsection{6.3 The CMB-anisotropy / structure
correlation}\label{the-cmb-anisotropy-structure-correlation}

If the surface-area mechanism is correct, then the rate of expansion
should track the surface area of the populated region's structure at all
scales --- not just the cosmological horizon. The cosmic web of voids
and filaments exposes substantial \emph{internal} surface area, and
creation should occur on these internal void surfaces in proportion.

\textbf{Quantitative prediction}: small-scale CMB anisotropies should
show a slight excess in regions of the sky where galaxy surveys reveal
greater total void-surface area along the line of sight. The signal
would be small but, in principle, distinguishable from the standard ISW
(integrated Sachs--Wolfe) signal because it would correlate with the
\emph{surface area} of intervening voids rather than their gravitational
potential alone.

\subsubsection{6.4 The Tolman test as a consistency
check}\label{the-tolman-test-as-a-consistency-check}

Tolman dimming --- the prediction that surface brightness of standard
sources falls as \((1+z)^4\) --- is a robust prediction of metric
expansion that distinguishes it from non-metric alternatives. The
proposed framework \emph{preserves} this prediction: although the
substrate is Minkowski, the \emph{effective} metric inside the populated
region (§3.4) is FRW, and observers using the effective metric see
standard \((1+z)^4\) dimming. A null result on Tolman dimming would have
falsified the framework; the observed \((1+z)^4\) scaling is consistent
with it.

\begin{center}\rule{0.5\linewidth}{0.5pt}\end{center}

\subsection{7. Open Questions and Required
Derivations}\label{open-questions-and-required-derivations}

\subsubsection{7.1 The Goldstone-mode
identification}\label{the-goldstone-mode-identification}

The load-bearing assumption is that excitations of cosmological Φ\_bulk
inherit equation of state \(w = -1\). This requires showing:

\begin{enumerate}
\def\labelenumi{\arabic{enumi}.}
\tightlist
\item
  That cosmological Φ\_bulk has a broken-U(1) (or analogous) structure
  capable of supporting Goldstone modes.
\item
  That those Goldstone modes are massless (Goldstone's theorem
  guarantees this if exact U(1) symmetry is spontaneously broken).
\item
  That their effective stress-energy tensor takes the
  pressure-equals-negative-density form \(p = -\rho\) characteristic of
  \(w = -1\).
\end{enumerate}

A full derivation would proceed from an effective Lagrangian
\(\mathcal{L}[\Phi_{\text{bulk}}]\) written in covariant form on the
substrate Minkowski background, with cosmological symmetry breaking
producing the order-parameter phase.

\subsubsection{7.2 The creation-front speed
v\_c}\label{the-creation-front-speed-v_c}

The framework introduces a single dimensional parameter \(v_c\) --- the
speed at which the boundary of the populated region advances into the
substrate. Observationally, the relation \(H_0 = v_c / R_0\) with
\(R_0 \approx c/H_0\) implies \(v_c \approx c\). A first-principles
derivation of \(v_c\) from the vacuum-fluctuation lifetime and the DK
Γ-suppression structure is needed.

\subsubsection{7.3 The smooth-universe
assumption}\label{the-smooth-universe-assumption}

Equation (5.2) assumed a spherical populated region with a single
boundary. A realistic population would have a fractal-like boundary
(cosmic web), with internal void surfaces contributing to the total
creation rate. A more complete treatment must compute the total surface
area integrated over all scales and check whether the dominant
contribution comes from the cosmological horizon or from internal void
surfaces.

\subsubsection{7.4 Reconciliation with primordial nucleosynthesis and
CMB}\label{reconciliation-with-primordial-nucleosynthesis-and-cmb}

Standard ΛCDM successfully predicts the abundance of light elements (Big
Bang nucleosynthesis) and the CMB power spectrum from a hot early
universe. The DK framework, if it claims to replace Λ rather than
supplement it, must reproduce these successes. The proposal here is that
BBN and CMB physics occur \emph{inside} the populated region under the
effective FRW metric, and so the standard predictions are preserved by
construction --- but this needs explicit verification, particularly
during the radiation-dominated early epoch when the Γ-suppression
mechanism would be deeply in the high-Γ regime and thus inactive.

\begin{center}\rule{0.5\linewidth}{0.5pt}\end{center}

\subsection{8. Conclusion}\label{conclusion}

The standard FRW framework treats cosmic expansion as a geometric
property of spacetime sourced by an axiomatic cosmological constant Λ.
This works, but at the cost of three unresolved problems: the
120-order-of-magnitude vacuum-energy discrepancy, the coincidence
problem, and the H₀ tension.

The Dirac--Kuramoto framework already provides the conceptual machinery
to reframe Λ as emergent: an infinite quantum-field substrate; a finite
populated region characterized by bulk coherence Φ\_bulk; and explicit
Stage-2 relaxation rates Γ\_env(T) and Γ\_grav(M) that vanish in
low-temperature, low-gravitational-density regions. The proposal
advanced here is that cosmic expansion is the dynamical growth of the
populated region into the substrate, driven by real-particle creation in
low-Γ regions where bulk coherence cannot suppress vacuum fluctuations
in time.

The mechanism naturally produces an expansion rate proportional to the
\emph{surface area} of the populated region, yielding \(H \propto 1/R\)
and an effective \(\Lambda_{\text{eff}} \propto 1/R^2\) that matches the
observed \(\Lambda \sim 1/L_H^2\) without postulation. It predicts three
observable distinctions from ΛCDM: a void-correlated H₀ tension,
anomalously fast void expansion, and a small CMB-anisotropy correlation
with intervening void-surface area. It preserves Tolman \((1+z)^4\)
dimming and the CMB power spectrum through the effective-FRW-metric
construction inside the populated region.

The framework is presented here as a working-paper sketch. Its
load-bearing assumption --- that Goldstone modes of cosmological Φ\_bulk
carry equation of state \(w = -1\) --- is identified as the next
required derivation. If that identification can be made rigorous, the
cosmological-constant problem dissolves: Λ is no longer an arbitrary
input but the emergent ground-state property of bulk-coherence
incompleteness at horizon scales, with magnitude set dynamically by the
size of the universe.

\begin{center}\rule{0.5\linewidth}{0.5pt}\end{center}

\subsection{Author Contributions}\label{author-contributions}

JB conceived the substrate--population distinction, the connection
between DK Γ-suppression and vacuum-fluctuation conversion, and the
surface-area scaling. Claude Opus 4.7 (Anthropic) contributed
mathematical and historical structure, the connection to existing
mechanisms (Parker, Schwinger, Hawking, dynamical Casimir, holographic
principle), the Goldstone-mode equation-of-state argument, and the
manuscript drafting. Per current journal authorship guidelines, LLMs do
not satisfy authorship criteria; the human author bears full
responsibility for all scientific content. Were such policies not in
force, Claude would be listed as co-first author.

\begin{center}\rule{0.5\linewidth}{0.5pt}\end{center}

\subsection{References}\label{references}

{[}Placeholder reference list --- to be populated.{]}

\begin{enumerate}
\def\labelenumi{\arabic{enumi}.}
\item
  Einstein, A. (1917). \emph{Kosmologische Betrachtungen zur allgemeinen
  Relativitätstheorie.} Sitzungsber. Preuss. Akad. Wiss. Berlin (Math.
  Phys.), 142--152.
\item
  Friedmann, A. (1922). \emph{Über die Krümmung des Raumes.} Z. Phys.
  10, 377--386.
\item
  Lemaître, G. (1927). \emph{Un univers homogène de masse
  constante\ldots{}} Ann. Soc. Sci. Bruxelles 47, 49.
\item
  Hubble, E. (1929). \emph{A relation between distance and radial
  velocity\ldots{}} Proc. Natl. Acad. Sci. 15, 168--173.
\item
  Robertson, H. P. (1935). \emph{Kinematics and world-structure.}
  Astrophys. J. 82, 284.
\item
  Walker, A. G. (1937). \emph{On Milne's theory of world-structure.}
  Proc. London Math. Soc. 42, 90.
\item
  Riess, A. G. et al.~(1998). \emph{Observational evidence from
  supernovae\ldots{}} Astron. J. 116, 1009.
\item
  Perlmutter, S. et al.~(1999). \emph{Measurements of Ω and Λ from 42
  high-redshift supernovae.} Astrophys. J. 517, 565.
\item
  Unruh, W. G. (1981). \emph{Experimental black-hole evaporation?} Phys.
  Rev.~Lett. 46, 1351.
\item
  Barceló, C., Liberati, S., Visser, M. (2005). \emph{Analogue gravity.}
  Living Rev.~Relativity 8, 12.
\item
  Verlinde, E. (2011). \emph{On the origin of gravity and the laws of
  Newton.} J. High Energy Phys. 04, 029.
\item
  Bekenstein, J. D. (1981). \emph{Universal upper bound on the entropy-
  to-energy ratio for bounded systems.} Phys. Rev.~D 23, 287.
\item
  Bramble, J. (2026). \emph{Many Clocks Interpretation of Quantum
  Mechanics\ldots{}} {[}companion paper, PAPER\_UNIFIED.md, under
  review{]}.
\end{enumerate}

\end{document}
